\chapter{18}

\section{Zürich (CH), Freitag, 24.
August}



Schliesslich war Valentina Giorgia hinterher gehetzt.
Die darauffolgende Stille im Raum war mächtig. Louis war hinter dem
Klavier zusammengesunken. Marco hatte sich als Nächsten bewegt. Er stand
auf, ging zu ihm und legte seine Arme von hinten über Louis’ Brust.

\sectionbreak

Tiefer hätte Louis nicht mehr fallen können. Giorgia musste wirklich
komplett abgeschlossen haben, mit ihm und mit ihrer Beziehung. Wieder
kamen dieselben Gedanken wie damals. Kapitulation, komplett versagt,
rettungsloser Idiot. Und diesmal öffentlich.
Seine Verzweiflung und sein Gesang in dieser Zerbrechlichkeit konnten ihr
nur noch den Rest gegeben haben. Was sollte sie mit einem wie ihm,
der nichts als bedürftig um Liebe bettelte. Wie ein jammerndes Kind. Als
unattraktiv musste sie ihn empfunden haben. Und dieser Brief, eine einzige Zumutung, für sie wie für ihn. 
Er hatte Tage damit verbracht, die Sätze zu formulieren. Hatte sich gequält und sich schliesslich durchgerungen,
das Kuvert wegzuschicken. Gegen all seine Bedenken. Sie musste seine
Worte als Erpressung empfunden haben.
Wie unfähig er immer wieder gewesen war. Wie schnell er hochkochte, 
sobald Beziehungen eine Grenze der Nähe überschritten, \emph{seine}
Grenze, machte er dicht. Und jetzt, als er seine innere Not offenbarte,
war es auch falsch.

\sectionbreak

Louis stach mit der Gabel das letzte Fruchtstück auf. Hatte er ihre Beziehung jemals
aus der Sicht von Giorgia betrachtet?
Er hatte versagt. Und war zum Feigling geworden.

\sectionbreak

Seine vergangene Nacht sollte in dieser Stadt die letzte gewesen sein. Versteckt,
in Daueranspannung, einmal mehr draussen, diesmal direkt am Fluss. Er
hatte überraschend gut geschlafen. Dennoch, es musste weg. Mit der einen
Hand hielt er die Kaffeetasse umschlossen, mit der andern winkte er den
Kellner zu sich.
Immerhin, er hatte es geschafft. In wenigen Tagen war er an einen neuen
Pass gekommen. Louis bezahlte.
Nach ziellosem Herumgehen stand er vor der Limmatbrücke. Er überquerte
sie und setzte sich am Quai auf die Mauer am Fluss. Immerhin, er war an die richtigen Menschen geraten. 
Sie hatten ihm geholfen. Ob das etwas mit Glück zu tun hatte? Vielleicht 
würde es unter diesem guten Stern weitergehen. Insgeheim glaubte er nur
vage daran.

\sectionbreak

\indentedblock{Damals. Fribourg, sechs.\\
Vorne geht Louis die Rolltreppe des Schuhgeschäftes empor. Er rennt
dazu.
«Warte», ruft Maman hinter ihm her, «kannst den Karton selbst tragen.»
Louis bleibt oben schliesslich stehen.
«Och nö,» ruft er zurück. Er hat sie ja, seine blauen Sneakers.
Nun ist auch Grandpère oben.
«Doch, Louis. Du hast genau die Schuhe bekommen, die du dir
wünschtest…und tragen kannst du sie selbst,» sagt Maman.
Sie hebt den Karton vor Louis’ Brust. Widerwillig hält er ihn und
trottet vorwärts.
Grandpère geht neben ihm. Louis schaut zu ihm hoch. Grandpère macht ein
ernstes, nachdenkliches Gesicht. Der Ausdruck ist Louis neu. So neu,
dass er erschrickt. Streng sieht Grandpère aus, fast böse, während er
eindringlich zu ihm sagt:
«Tu sais, mon cher\footnote{«Weisst du, mein Lieber,…»}, für Geschenke bedankt man sich. Ein Danke kostet
dich nichts.»
Der soll sich nicht einmischen, denkt Louis und vergrössert den Abstand
zu Grandpère. Sein Vater will, dass er die schönsten Dinge hat, denkt er
trotzig. Je teurer, desto besser. Er kauft ihm alles, was er will. Und
Schuhe braucht er so oder so. Louis sieht unsicher in Mamans Gesicht.
Ist sie traurig?
Jetzt schaut Grandpère wieder ein bisschen freundlicher. Louis atmet
auf. Hat er vielleicht doch recht? Louis überlegt.
«Merci, Maman,» sagt er schnell und leise.
Maman nickt sanft und umarmt ihn. Es fühlt sich gut an.
Sie setzen sich zusammen in das kleine Café. Louis bekommt eine heisse
Schokolade.
Dann wechselt er die Schuhe. Die alten Sandalen legt er in die
Schachtel.
«Sie sind schön, Maman,» sagt er und blickt an sich hinunter. Dann
strahlt er die beiden abwechselnd an. Grandpère klopft ihm sanft auf die
Schulter.}

\sectionbreak

Erst viele Jahre später hatte es Louis verstanden. Maman hatte ums Geld
gekämpft. Um jede Münze. Es hatte schon weit vor Vaters Verschwinden begonnen. Später
waren Unterhaltszahlungen vereinbart, um konsequent auszubleiben. Vaters
Insolvenz folgte kurz nach der Scheidung. Unlautere Machenschaften
flogen auf. Die Autogarage wurde geschlossen. Und immer wieder
verschwand Vater aus unerklärlichen Gründen nach Spanien, sagte er
jedenfalls. Grandpère half, so gut es ging. Und Maman hatte Lucia und
ihn vor den Tatsachen verschont. Über all die Jahre.
«Du hast Papa so unendlich verehrt, Louis,» hatte sie viel später zu ihm
gesagt. An einem Abend, als er endlich alles wissen wollte.
«Du warst zu klein, Lucia noch kleiner, ihr wart zu jung, um zu
verstehen. Und ihr hattet das Recht auf eine Kindheit mit beiden
Eltern.»
Louis’ angekratztes Vaterbild war längst in Scherben zerfallen. Er hatte
Monate gebraucht, seine Wut zu verstehen. Auch die auf Maman.

\sectionbreak

Seine letzte Handlung in Zürich sollte der Kauf einer Gitarre sein.
Endlich. Als Ciro, sozusagen. Und Louis sah sich und sein Instrument.
Seine «Taylor». Er hörte die Klänge spielen, versank summend darin und
blickte ins Nichts. 

Er stand auf und setzte sich in Bewegung. Nach wie vor ahnungslos,
wohin ihn der Weg führen sollte. Doch ihm fiel auf, dass er geschmeidiger
und aufrechter ging. Aber auch, wie genau er die Umgebung beobachtete
und dauernd die Furcht wieder hochkam, aufgegriffen zu werden. Als würde
sie ihn im Würgegriff halten, kurz pausieren, um gleich resolut nachzufassen.
Das nächstliegende Gitarrengeschäft fand er am \emph{Limmatquai}. Auf
schwarzem Grund mit weisser Schrift las er «Musik Hug» über dem Eingang.
Louis erkannte bereits durch die Schaufenster, dass ihn das Angebot
nicht interessierte.
Er schleppte sich weiter. Über den Hauptbahnhof gelangte er an die
\emph{Langstrasse}. In der Hoffnung auf einen schattigen Platz liess er
sich treiben. Nach über einer Stunde schwand sein Antrieb. Er nervte sich über sich selbst. Was machte er hier eigentlich?
Er entschied, ein Restaurant zu suchen. An einer Tramstation blieb er
stehen. Wie es aussah, befand er sich im Gebiet \emph{Aussersihl}.
Während er weiterging, glaubte er in der Ferne eine kleine, bunte Oase
zu entdecken. Es deutete sich eine Parkanlage an.
Je näher er kam, desto klarer konnte er die Beschriftung eines grünen,
modernen Gebäudes erkennen. «Quartierzentrum - Restaurant Bäckeranlage».
Das Haus lag vor einer ausgedehnten Rasenfläche. Darauf schienen sich
uralte, stattliche Bäume in Szene zu setzen. Sie verbreiteten eine
wohlwollende Stimmung.
Die Tische vor dem Restaurant waren alle besetzt.
Louis durchschritt den Park.
Am Rande der Wiese sah er einen Mann im Schatten unter einer mächtigen
Rotbuche liegen. Ihr sattes Dunkelgrün schien der Trockenheit zu
trotzen.
Der Mann wirkte sympathisch. Die Arme hatte er unter den Kopf gelegt.
An den Stamm gelehnt stand eine Gitarre. Und Louis zielte auf geradem
Weg auf ihn zu.
«Hey, ich bin Ciro, darf ich?»
Louis deutete mit der Hand auf den Boden neben dem Mann.
Der schien seine Worte nicht zu verstehen, Louis’ Geste schon. Er
lächelte.
«Of course, please,» sagte er und setzte sich auf.
Er war Engländer und stellte sich als Dave vor.
Ihr Gespräch wurde im Nu zu einem angeregten Austausch über Musik,
Interpreten und Stilrichtungen.
Sie hatten sich in kurzer Zeit gefunden und blieben bei «Tom Waits»
hängen. Daves Gesicht leuchtete.
«Kennst du «Make it rain» von ihm oder \emph{«Hold on»}?»
Dave schaute Louis fragend an. Louis nickte.
«Ja, sehr schön. «Hold on» gefällt mir gut, auch den Text mag
ich, so was von Durchhalten und Hoffnung, nicht?»
«Yes.»
Dave schaute zu Louis hoch, als fühle er sich verstanden.
«Mags du’s spielen?»
Dave wirkte, als hätte er auf die Frage gewartet. Und schon packte er
seine Gitarre, setzte sich bequem hin und senkte die Augen.
Louis war auf der Stelle gefangen. Daves subtil-virtuose Spielweise war
der Ausdruck eines Könners. Die leicht verzögerten Klänge legten sich
breit unter den Baum. Die Sommerwiesenstimmung kippte in süsse Schwere.
Daves Stimme, rau und verblüffend tief, fesselte Louis. Die Worte trafen
ihn, trugen ihn davon, um ihn augenblicklich wieder zurückzuschleudern.
\emph{Du musst durchhalten. Du musst durchhalten.} Louis war, als reiche
ihm Dave gerade die Hand.
Er hörte ehrfürchtig zu. Weder die Hitze noch die sich rundum nähernden
Menschen nahm er mehr wahr.

\qrblock{Tom Waits \emph{«Hold on»}}{media/QR-Codes/018_Hold_on_Tom_Waits_Spotify.png}

Dave spielte den letzten Akkord. Erst blieb es still.
Dann setzte anerkennendes Klatschen ein. Einer rief begeistert «Genial!»
Die meisten blieben stehen und schauten Dave bewundernd an.
Louis beugte sich zu ihm hinüber und klopfte ihm auf die Schulter. Seine
Augen waren glasig. Dave war sichtlich gerührt. Die Menschen zerstreuten
sich.
Louis war wieder mit ihm alleine.
